\begin{figure}[htbp]
\centering
\begin{tikzpicture}[x=1.05cm,y=1.05cm,
  every node/.style={font=\small},
  source/.style={thick,-{Stealth[length=2mm]}},
  charge/.style={thick,dashed,-{Stealth[length=2mm]}}]
\fill[gray!14] (0,1.3) -- (2,1.3) -- (2,0)
  -- (2.6,0) -- (2.6,1.9) -- (0,1.9) -- cycle;
\draw[very thick] (-2,-1.3) rectangle (2,1.3);
\node at (-.65,-.35) {$\Omega$};
\node at (3.3,-.85) {$O$};
\node at (-1.55,.8) {$A$};
\node at (-1.55,-.85) {$B$};
\node at (1.55,-.85) {$C$};
\node at (1.55,.85) {$D$};
\foreach \point in {(0,1.3),(-2,0),(0,-1.3),(2,0)}
  \fill \point circle[radius=1.5pt];
\node[below left] at (0,1.3) {$p$};
\node[left] at (-2,0) {$b$};
\node[below] at (0,-1.3) {$c$};
\node[below right] at (2,0) {$d_0$};
\draw[source] (0,1.3) .. controls (.1,.95) and (.25,.65) .. (.25,.25);
\fill (.25,.25) circle[radius=1.5pt];
\node[right] at (.25,.25) {$z$};
\draw[source] (0,1.3) .. controls (.35,1.73) and (1.75,1.7) .. (2.28,1.6)
  .. controls (2.4,1.45) and (2.4,.95) .. (2.4,.65);
\fill (2.4,.65) circle[radius=1.5pt];
\node[right] at (2.4,.65) {$u$};
\node[above,align=center] at (1.45,1.93) {exterior $D$ strip};
\draw[charge] (0,1.3) .. controls (-.4,1.83) and (-1.85,2.05) .. (-2.65,2.35);
\draw[charge] (-2.65,2.35) .. controls (-3.05,2.52) and (-3.45,2.9) .. (-3.75,3.1);
\draw[charge] (-2.65,2.35) .. controls (-3.05,2.25) and (-3.45,2.15) .. (-3.75,2.05);
\fill (-3.75,3.1) circle[radius=1.5pt];
\fill (-3.75,2.05) circle[radius=1.5pt];
\node[left,align=right] at (-3.8,3.1)
  {$w_j^{(1)}$\\$\sum_j a_j^{(1)}=1$};
\node[left,align=right] at (-3.8,2.05)
  {$w_j^{(2)}$\\$\sum_j a_j^{(2)}=1$};
\node[above] at (-1.65,2.3) {charge $-2$};
\node[align=center] at (-3.05,.65) {outgoing cut\\on the $A$ side};
\node[align=center] at (3.6,-.05)
  {opposite wire\\colors on\\the two banks};
\end{tikzpicture}
\caption{The paired-wall geometry in the original plane.  The solid
arrows are the two source paths from $p$; the exterior probe stays in
the shaded strip along $D$.  The dashed pure strings leave on the
other side of $p$.  Each of their two endpoint lists has total charge
$-1$, so their common initial cut has charge $-2$.  Only representative
distant endpoints are shown, not their separated metric scales.
The polygon is schematic; the short interval strings and the fine
parallel lanes at $p$ are omitted.  Inversion is used only to describe
the ideal exterior turn factors, not to change its probability law.}
\label{wall:geometry-figure}
\end{figure}
